\chapter{Elementary modules and the PID structure theorem}\label{chap:M07.16}

% Topic grouping: The common language
\KEntry{K07.16.definition}{Modules}
A unital left $R$-module is an abelian group $M$ with scalar action
such that
\[
 1m=m,\quad (rs)m=r(sm),\quad
 (r+s)m=rm+sm,\quad r(m+n)=rm+rn .
\]
These axioms imply $0m=0$, $r0=0$, $(-r)m=-(rm)$ and $r(-m)=-(rm)$
by the additive cancellation arguments used for rings.
Vector spaces are modules over fields. Every abelian group has the
$\Z$-module structure given by repeated addition and negatives; the
integer distributive and associative laws follow by counting finite sums
and then using additive inverses. Conversely a $\Z$-module has precisely
this action, since $1m=m$ forces positive multiples and negatives.
Left ideals are submodules of the left regular module $R$.
Modules have no assumed multiplication of their own elements; left and
right scalar conventions must not be confused over a noncommutative ring.

\KEntry{K07.16.maps_submodules}{Linear maps and submodules}
An $R$-linear map preserves addition and scalar multiplication.
A submodule is an additive subgroup closed under scalar multiplication.
Equivalently it is a nonempty subset closed under addition and scalar
multiplication: the latter conditions give $0=0m$ and $-m=(-1)m$.
Kernels and images of linear maps are submodules, since
$f(x+y)=f(x)+f(y)$ and $f(rx)=rf(x)$ preserve zero and describe images
of sums and scalar multiples. Intersections are submodules by testing
each condition in each member. The span of a set $S$ consists of finite
sums $\sum r_i s_i$; it is closed under these operations, contains $S$,
and lies in every submodule containing $S$. This proves its universal
least-submodule property without assuming a basis.

\KEntry{K07.16.quotients_isomorphism}{Quotients and cyclic modules}
For $N\le M$, define $r(m+N)=rm+N$. Replacing $m$ by $m+n$
changes the result by $rn\in N$, so the action is well-defined; the
module identities descend from $M$. If $f:M\to L$ is linear and
$N\subseteq\ker f$, the formula $\bar f(m+N)=f(m)$ is well-defined
and linear, and is the unique map with $\bar f\pi=f$ because $\pi$ is
onto. Taking $N=\ker f$ and target $\im f$ gives a bijection: equal
values mean the representatives differ by a kernel element. Thus
\begin{equation}\label{eq:module-first-isomorphism}
 M/\ker f\ \cong_R\ \im f .
\end{equation}
If $R$ is commutative, $M=Rm$, and $I=\Ann_R(m)=\{r:rm=0\}$,
the onto linear map $R\to M$, $r\mapsto rm$, has kernel $I$.
Consequently $R/I\cong_R M$ by $r+I\mapsto rm$.
This is an isomorphism of modules; no ring multiplication on $M$ is
asserted. Over a noncommutative ring the annihilator of a generator
can be only a left ideal, and the module quotient need not be a quotient
ring. For example, in the standard $M_2(F)$-module $F^2$, the
annihilator of $e_1$ consists of matrices with first column zero.
It contains $E_{12}$ but not $E_{12}E_{21}=E_{11}$.

\KEntry{K07.16.free_generators}{Bases, finite generators, and rank}
A basis $(e_i)$ expresses each element uniquely as a finite sum
$\sum r_i e_i$; a module admitting a basis is free.
Finite generation requires a finite spanning set, which need not be
independent. The module $\Z/2\Z$ is generated by one element but is
not free: a nonzero free $\Z$-module has no nonzero element annihilated
by $2$, by uniqueness of basis coefficients.

For a domain $R$, finite free rank is invariant. Indeed an isomorphism
$R^m\cong_R R^n$ is represented by matrices $A,B$ with $AB=I_n$,
$BA=I_m$. Embed all entries into $\Frac(R)$; the same identities give
inverse linear maps between vector spaces of dimensions $m,n$ over that
field. The dimension theorem from D05 gives $m=n$. This proof covers
rank zero and uses only finite-dimensional vector-space theory.

\begin{lemma}\label{lem:pid-primitive-basis}
Let $R$ be a PID. A vector $e\in R^n$ whose coordinates generate $R$
extends to a basis.
\end{lemma}
\begin{proof}
For coordinates $x,y$ not both zero, choose a generator $d$ of $(x,y)$.
Write $x=dx'$, $y=dy'$ and $d=rx+sy$. Cancellation gives
$rx'+sy'=1$. The matrix
\[
 \begin{pmatrix}r&s\\-y'&x'\end{pmatrix}
\]
has determinant $1$, inverse
$\begin{pmatrix}x'&-s\\y'&r\end{pmatrix}$, and sends
$(x,y)^t$ to $(d,0)^t$. A pair of zero coordinates needs no operation.
Apply such invertible changes successively to the first nonzero
coordinate and each remaining coordinate. The resulting vector is
$(g,0,\ldots,0)$, where $(g)$ is the ideal generated by all original
coordinates. Since it is $R$, $g$ is a unit; scale that coordinate by
$g^{-1}$. The resulting invertible transformation sends $e$ to the
first coordinate vector, so its inverse transports the coordinate basis
to a basis beginning with $e$.
\end{proof}

% Topic grouping: A diagonal basis over an arbitrary PID
\KEntry{K07.16.submodule_free}{Submodules of finite free modules}
\begin{theorem}\label{thm:pid-smith-diagonalization}\label{thm:pid-submodule-free}
Let $M$ be free of finite rank $n$ over a PID $R$ and $N\le M$.
There are a basis $e_1,\ldots,e_n$ of $M$ and nonzero elements
$a_1,\ldots,a_s$, $s\le n$, with
$a_1\mid a_2\mid\cdots\mid a_s$, such that
$a_1e_1,\ldots,a_se_s$ is a basis of $N$.
In particular every submodule of a finite free module is finite free.
\end{theorem}
\begin{proof}
Induct on $n$. If $n=0$ or $N=0$, use the empty basis of $N$.
Otherwise consider the ideals
$\{\varphi(N):\varphi\in\Hom_R(M,R)\}$.
These are ideals because $\varphi$ is linear. By
Lemma~\ref{lem:pid-acc} this nonempty collection has a maximal member,
say $\varphi(N)=(a)$. Some coordinate projection is nonzero on $N$,
so zero cannot be maximal; hence $a\ne0$.
Choose $y\in N$ with $\varphi(y)=a$.

We first prove that $a$ divides $\psi(y)$ for every linear functional
$\psi:M\to R$. Let $(d)=(a,\psi(y))$ and
$d=ra+s\psi(y)$, and put $\theta=r\varphi+s\psi$.
Then $d=\theta(y)$, so
$(a)\subseteq(d)\subseteq\theta(N)$.
The last ideal belongs to the chosen collection, whose maximal member
is $(a)$. Thus $(d)=(a)$, proving the divisibility.
Apply this to all coordinate projections: all coordinates of $y$ are
multiples of $a$, so $y=ae$ for some $e\in M$.
The equation $a=\varphi(y)=a\varphi(e)$ gives $\varphi(e)=1$.
Hence the coordinates of $e$ generate $R$, since $1=\varphi(e)$
is a linear combination of them. Lemma~\ref{lem:pid-primitive-basis}
extends $e$ to a basis $e,u_2,\ldots,u_n$ of $M$.

Set $M_0=\ker\varphi$. The vectors
$u_i-\varphi(u_i)e$ for $i\ge2$ form a basis of $M_0$.
Indeed subtracting $\varphi(m)e$ from a basis expansion of $m$
shows spanning, and comparing the original $u_i$ coefficients shows
independence. Thus $M=Re\oplus M_0$ with $M_0$ free of rank $n-1$.
For $x\in N$, write $\varphi(x)=ac$. Then $x-cy\in N\cap M_0$,
and the intersection of $Ry$ with $M_0$ is zero because $a\ne0$.
Consequently
\[
 N=R(ae)\oplus N_0,\qquad N_0=N\cap M_0.
\]

We also need the divisibility joining this first summand to the rest.
If $\psi$ vanishes on $e$, then
$(\varphi+\psi)(y)=a$. Maximality again gives
$(\varphi+\psi)(N)=(a)$. For $x\in N_0$, this says $\psi(x)\in(a)$.
Use the coordinate functionals of the displayed basis of $M_0$, extended
by zero on $e$. Every coordinate of every $x\in N_0$ is divisible by
$a$, so $N_0\subseteq aM_0$.

Apply induction to $N_0\le M_0$. It gives a basis $e_2,\ldots,e_n$
of $M_0$ and a basis $a_2e_2,\ldots,a_se_s$ of $N_0$, with a
divisibility chain among these nonzero coefficients.
Since $a_i e_i\in aM_0$, its $e_i$ coefficient is divisible by $a$.
Thus $a\mid a_i$ for every $i$, in particular $a\mid a_2$ if present.
Taking $e_1=e$, $a_1=a$ gives the required complete chain and basis.
No Euclidean division or classification of abelian groups was used.
\end{proof}

\begin{corollary}\label{cor:pid-submodules-finite}
Every submodule of a finitely generated module over a PID is finitely
generated.
\end{corollary}
\begin{proof}
Choose an onto linear map $\pi:R^n\to M$ from a finite generating set.
For $L\le M$, the inverse image $\pi^{-1}(L)$ is a submodule of $R^n$,
hence has a finite basis by the theorem. The images of its basis vectors
span $L$: each $l\in L$ has a preimage in that inverse image.
\end{proof}

\KEntry{K07.16.finite_presentation}{Finite presentations}
A finite presentation is an onto map $\pi:R^n\to M$ whose kernel is
finitely generated. Choosing kernel generators gives a finite matrix
$A:R^m\to R^n$ with $\im A=\ker\pi$, and hence
$M\cong_R R^n/\im A=\coker A$ by
\eqref{eq:module-first-isomorphism}. Conversely such a finite matrix
has its finitely many column vectors as image generators, so defines a
finite presentation. Over a PID every finitely generated module has
one, because the kernel of its chosen $R^n$ surjection is finite free.

Finite generation alone is insufficient over a general ring.
Let $R=k[x_1,x_2,\ldots]$ be the union of its finite-variable polynomial
rings and $I=(x_1,x_2,\ldots)$. The module $R/I\cong k$ is cyclic,
but $I$ is not finitely generated. In fact any finite list in $I$
involves only finitely many variables. Choose $x_j$ outside that list;
setting all other variables to zero maps every listed generator to zero
but $x_j$ to a nonzero indeterminate. Thus the list cannot generate $I$.

We justify why a different finite presentation cannot evade this
obstruction. Suppose $\pi:R^n\to R/I$ is onto with finitely generated
kernel $L$. Let $q:R\to R/I$ be the natural map.
Choose linear lifts $u:R\to R^n$ of $q$ and $v:R^n\to R$ of $\pi$,
by choosing a preimage of $q(1)$ and of each $\pi(e_i)$.
Then $\pi u=q$ and $qv=\pi$.
For $x\in I$, $u(x)\in L$ and
$x=(1-vu)(x)+vu(x)$.
The image of $1-vu$ is generated by the one element $(1-vu)(1)$
and lies in $I$; also $v(L)\subseteq I$.
Hence $I$ is generated by that one element and the images of a finite
generating set of $L$, contradicting the preceding paragraph.

\KEntry{K07.16.torsion}{Torsion and the torsion-free quotient}
For a module over a domain, call $x$ torsion when some $r\ne0$ kills
it. Denote its torsion submodule by $T(M)$.
If $rx=0$, $sy=0$ with $r,s\ne0$, then $rs(x+y)=0$ and $rs\ne0$.
Negatives and scalar multiples remain torsion by commutativity.
Thus $T(M)$ is a submodule. The quotient is torsion-free: if
$r(x+T)=0$, $r\ne0$, then $rx\in T$ and some $s\ne0$ gives
$srx=0$, so $x\in T$.
For a finitely generated PID-module, $T(M)$ has finite generators
$t_1,\ldots,t_l$ by Corollary~\ref{cor:pid-submodules-finite}.
Choose $r_i\ne0$ killing $t_i$; the nonzero product $\prod r_i$
kills all of $T(M)$. For $T=0$ use the scalar $1$.
Torsion need not mean finite cardinality: $k[X]/(X)$ is torsion over
$k[X]$ and is infinite if $k$ is infinite.
The quotient by torsion is canonical, but a complementary free submodule
need not be. In $\Z\oplus\Z/2\Z$ the subgroups generated by $(1,0)$
and by $(1,[1])$ are distinct. Each meets the torsion subgroup
$0\oplus\Z/2\Z$ trivially and, together with it, generates the whole
module, by subtracting the appropriate integer multiple of its generator.

\KEntry{K07.16.smith}{Maps, Smith form, and cancellation of presentations}
\begin{corollary}\label{cor:pid-presentation-cokernel}
A map $A:R^m\to R^n$ over a PID admits bases in which it has a
diagonal block with nonzero entries $a_1\mid\cdots\mid a_s$ and
all other entries zero. Its cokernel is
\[
 \coker A\ \cong_R\ R^{n-s}\oplus\bigoplus_{i=1}^s R/(a_i).
\]
Unit $a_i$ give zero quotient summands and may be cancelled from the
presentation. Zero or unused target rows give free summands.
\end{corollary}
\begin{proof}
Apply Theorem~\ref{thm:pid-smith-diagonalization} to
$N=\im A\le R^n$, obtaining a target basis $e_i$ and image basis
$a_i e_i$. Choose $u_i\in R^m$ with $A(u_i)=a_i e_i$.
Every domain vector differs from a linear combination of the $u_i$
by an element of $\ker A$, and their span intersects $\ker A$ trivially,
since the $a_i e_i$ are independent.
The kernel is finite free by the same submodule theorem.
Joining a kernel basis to $u_1,\ldots,u_s$ therefore gives a domain
basis. Its length is $m$ by rank invariance.
The resulting matrix is the asserted diagonal matrix.
Taking target coordinates modulo $a_i$ in the first $s$ positions and
leaving the rest unchanged defines an onto map to the displayed direct
sum, with kernel precisely $N$. The first isomorphism theorem proves
the cokernel formula.
Finally $R/(a_i)=0$ for a unit, while an unused target coordinate is
unrestricted, so deleting unit relation/generator pairs preserves the
cokernel.
\end{proof}
This last deletion changes presentation ranks. It is not an invertible
basis change of the original fixed-rank map: the identity $\Z\to\Z$
still has a unit matrix at rank one, whereas its zero cokernel also has
the empty presentation.

\begin{corollary}\label{cor:integer-presentation-cokernel}
For an integer $n$-by-$m$ presentation matrix, its cokernel is
$\Z^{n-s}\oplus\bigoplus_i\Z/d_i\Z$, with positive
$d_1\mid\cdots\mid d_s$ after changing signs. Factors $d_i=1$ may
be deleted. If $s=n$, its finite cardinality is $\prod_i d_i$;
if $s<n$, it is infinite.
\end{corollary}
\begin{proof}
Apply the preceding corollary to the PID $\Z$, normalize nonzero
generators by their signs, and count independent coordinates modulo each
$d_i$. A nonzero free coordinate contains an injective copy of $\Z$.
\end{proof}

For a Euclidean domain one can compute the diagonalization by finite
row and column operations. Move a nonzero entry to the first position.
Divide first-column entries by that pivot and subtract row multiples;
a nonzero remainder can be swapped into the pivot position and strictly
decreases its Euclidean size. Do the analogous operation on the first
row using columns. If no nonzero remainder occurs, these operations
clear the row and column. If the pivot $a$ fails to divide an entry $b$
in the remaining block, add the row of $b$ to the first row. Its first
entry is still $a$, and its $b$ entry now yields a nonzero smaller
remainder by a column division; swap that remainder into the pivot.
Every restart thus decreases the pivot, so there are only finitely many
restarts; between restarts only finitely many row and column divisions
are made. Eventually the pivot divides the whole remaining block.
Recurse on that block. All subsequent operations preserve divisibility
by the first pivot, producing the divisibility chain. This is an
algorithmic supplement to, not an assumption in, the generic PID proof.

% Topic grouping: Existence and normalized uniqueness
\KEntry{K07.16.structure_existence}{The PID structure theorem}
\begin{theorem}\label{thm:pid-structure-existence}
Every finitely generated module $M$ over a PID has decompositions
\[
 M\cong_R R^r\oplus\bigoplus_{i=1}^s R/(d_i)
 \cong_R R^r\oplus\bigoplus_{p,j}R/(p^{e_{p,j}}),
\]
where $r,s\ge0$, the $d_i$ are nonzero nonunits with
$d_1\mid\cdots\mid d_s$, and the $e_{p,j}$ are positive integers,
only finitely many being present. Primes are taken up to associates.
The free part vanishes exactly when $M$ is torsion.
\end{theorem}
\begin{proof}
Choose a finite presentation and apply
Corollary~\ref{cor:pid-presentation-cokernel}. Delete unit entries, which
are an initial segment of the divisibility chain, and call the remaining
entries $d_i$. This gives the first decomposition.
Factor each $d_i$ into prime powers using
Theorem~\ref{thm:pid-implies-ufd}.
Distinct prime powers are comaximal: a gcd of them is a unit, and the
PID Bezout identity supplies their sum $1$.
The finite CRT therefore gives
$R/(d_i)\cong_R\bigoplus_{p\mid d_i}R/(p^{v_p(d_i)})$.
Here the ring CRT is also $R$-linear because all its coordinate maps
respect multiplication by base scalars. This gives the second form.
Each nonzero cyclic quotient is torsion; their finite direct sum is
annihilated by the product of their annihilators. A free nonzero module
over a domain has no nonzero torsion element, by coefficient comparison.
Thus $M$ is torsion exactly when $r=0$.
\end{proof}
As another consequence, a finitely generated PID-module is free if and
only if it is torsion-free. The forward implication is coefficient
cancellation. For the reverse, the displayed cyclic nonunit quotients
would contain nonzero torsion classes $[1]$; hence none can occur.
The finite-generation hypothesis cannot be dropped: $\Q$ is
torsion-free as a $\Z$-module but is not free. If a basis contained
$e\ne0$, write $e=2q$ and expand $q$ in that basis. The coefficient
of $e$ would give $1=2n$ with $n\in\Z$, a contradiction. An empty
basis cannot span the nonzero module $\Q$.
The PID hypothesis for submodule freeness also matters.
The ideal $(2,X)\subset\Z[X]$ is not free. Two nonzero elements
$a,b$ of an ideal in a commutative domain are dependent as module
vectors, because $b\cdot a-a\cdot b=0$ has nonzero coefficients.
Thus a nonzero free ideal would have a one-element basis and be
principal, contrary to the proof for $(2,X)$ in
Chapter~\ref{chap:M07.10}.

The free term must not be omitted. In Aluffi, Theorem 9.40 has it;
Corollary 9.44 as printed describes all finitely generated modules by
prime-power torsion summands alone. Taking $R=\Z$, $M=\Z$ disproves
that wording. The formula above is the corrected free-plus-torsion
statement. No classification of finite abelian groups was used in its proof.

\KEntry{K07.16.structure_uniqueness}{Uniqueness of both normal forms}
\begin{theorem}\label{thm:pid-structure-uniqueness}
In Theorem~\ref{thm:pid-structure-existence}, $r$ and the multiset of
prime powers $p^{e_{p,j}}$ are uniquely determined up to associates.
The ordered invariant factors $d_1,\ldots,d_s$ are consequently unique
up to associates, after excluding unit factors.
\end{theorem}
\begin{proof}
An isomorphism preserves torsion because it commutes with multiplication
by every scalar. For either decomposition its torsion submodule is
exactly the cyclic quotient part, so $M/T(M)\cong R^r$.
The rank invariance proved using $\Frac(R)$ gives uniqueness of $r$.
This is the generic, or fraction-field, rank; it does not count torsion.

For each irreducible $p$ define the intrinsic primary submodule
\[
 T_p=\{x\in T(M):p^k x=0\text{ for some }k\ge0\}.
\]
The sum and scalar-multiple tests show this is a submodule.
In a summand $R/(q^e)$ with $q$ not associate to $p$, Bezout gives
$up+vq^e=1$, so multiplication by $p$ is invertible on that summand.
It has no nonzero element killed by a power of $p$.
In a $p$-power summand every element is killed by that power.
Thus $T_p$ is exactly the direct sum of the $p$-power summands in
any primary decomposition.

The quotient
$p^{k-1}T_p/p^kT_p$, $k\ge1$, is a vector space over the field
$R/(p)$: the scalar action factors through that field because $p$
kills the quotient. Its dimension $b_{p,k}$ is an isomorphism invariant.
For a single summand $A=R/(p^e)$, this quotient is zero if $k>e$.
If $1\le k\le e$, the map
$R\to p^{k-1}A/p^k A$, $r\mapsto[p^{k-1}r]$, is onto and has
kernel $(p)$. Indeed membership in the kernel means
$p^{k-1}r=p^k s+p^e t$ in $R$ for some $s,t$; cancelling the nonzero
$p^{k-1}$ gives $r=ps+p^{e-k+1}t$, equivalent to $p\mid r$.
Hence the quotient is one-dimensional over $R/(p)$.
Finite direct sums commute with this filtration by taking coordinates,
so
\[
 b_{p,k}=\#\{j:e_{p,j}\ge k\},\qquad
 b_{p,k}-b_{p,k+1}=\#\{j:e_{p,j}=k\}.
\]
All exponents and their multiplicities are therefore intrinsic.

It remains to justify the ordered invariant-factor normalization.
For each prime, arrange its positive exponents increasingly. Let $s$
be the largest number of factors for any one prime, and pad every
shorter list on the left by zeros to length $s$.
If there are no prime factors, set $s=0$.
Multiply the prime powers columnwise, obtaining $d_1,\ldots,d_s$.
Each exponent is nondecreasing, so $d_i\mid d_{i+1}$.
CRT splits each column back into its nontrivial prime-power factors,
so this reconstructs the given torsion module.
Conversely, in a chain of nonunit invariant factors the exponent list
for each prime is nondecreasing. A prime dividing $d_1$ divides all
$s$ terms, so the largest count of positive exponents is exactly $s$.
The chain's exponent lists are consequently exactly the padded lists
just described. Unique factorization in $R$ determines each column
up to a unit, proving uniqueness.
If the torsion module is zero, all lists and the invariant-factor list
are empty and the argument still applies.
\end{proof}
Arbitrary coprime grouping would not give this normalization:
elementary divisors $2,4,3$ can be grouped as $6,4$, which have no
divisibility order; sorting and padding gives the invariant factors
$2,12$. This normalization is additional mathematics beyond merely
diagonalizing a matrix.

\KEntry{K07.16.operator_bridge}{Polynomial modules and canonical forms}
Let $V$ be finite-dimensional over $F$ and $T\in\End_F(V)$.
Define $f\cdot v=f(T)v$. Polynomial evaluation is a ring homomorphism
because scalars commute with $T$, so the module axioms hold. Conversely
an $F[X]$-module structure extending the given scalar action determines
$T(v)=Xv$, and distributivity shows that $T$ is $F$-linear and that
the original action is exactly evaluation at $T$.

An $F$-basis spans over $F[X]$ too, so this module is finitely generated.
An algebraic operator alone would not imply finite dimension: the zero
operator on a space with an infinite basis is killed by $X$.
It is torsion: for each $v$, the $d+1$ vectors
$v,Tv,\ldots,T^d v$ are dependent when $d=\dim_F V$, giving a
nonzero polynomial killing $v$. The structure theorem over the PID
$F[X]$ therefore gives
\[
 V\cong_{F[X]}\bigoplus_i F[X]/(f_i),\qquad
 f_1\mid\cdots\mid f_s,
\]
with monic nonconstant $f_i$. There is no free part, because a copy of
$F[X]$ has the infinite independent sequence $1,X,X^2,\ldots$ over $F$.
For $f=X^n+c_{n-1}X^{n-1}+\cdots+c_0$, the quotient basis
$1,X,\ldots,X^{n-1}$ makes multiplication by $X$ its companion matrix:
the first $n-1$ basis vectors shift forward and the last maps to
$-\sum c_iX^i$. Expansion of $\det(tI-C_f)$ down its last column gives
$t^n+\sum_{i=0}^{n-1}c_i t^i$: the cofactor corresponding to $c_i$
has $i$ diagonal $t$ factors, the remaining off-diagonal $-1$ factors
cancel its cofactor sign, and the final diagonal term gives $t^n$.
Thus each block has characteristic polynomial $f_i$.
Its annihilator as a polynomial module is $(f_i)$, because a polynomial
kills the generator $1$ exactly when it is zero modulo $f_i$.
On the direct sum the annihilator is their intersection
$(f_s)$, whereas block determinants multiply. Consequently
\[
 \operatorname{minpoly}(T)=f_s,\qquad
 \operatorname{charpoly}(T)=\prod_i f_i
\]
for nonzero $V$; on $V=0$ both polynomials are $1$ by the empty-block
and whole-annihilator conventions. Module isomorphisms are exactly
invertible linear maps commuting with $T$, since commuting with $X$
implies commuting with every polynomial. This proves the rational
canonical classification mechanism and its uniqueness.

If the $f_i$ split over $F$, the primary summands are
$F[X]/((X-\lambda)^e)$. In the reversed basis
$(X-\lambda)^{e-1},\ldots,1$, multiplication by $X$ is $\lambda\id$
plus the matrix with ones immediately above the diagonal: a Jordan
block. Thus the same theorem gives Jordan blocks precisely when the
needed factors split. The D05 results about matrices, change of basis,
determinants and Cayley--Hamilton keep their existing home; this argument
provides the module representation and classification bridge rather
than another independent list of operator invariants.
